\section*{Введение}
\addcontentsline{toc}{section}{Введение}

Теория графов~--- раздел дискретной математики, изучающий объекты и связи между ними с помощью графов, состоящих из вершин и рёбер. В рамках теории графов рассматриваются задачи анализа связности, поиска путей, построения остовов, исследования циклов, потоков и паросочетаний. Графовые модели применяются при описании транспортных и компьютерных сетей, социальных связей, структур данных и других сложных систем.

В отчёте представлено описание выполнения цикла лабораторных работ по дисциплине <<Теория графов>>. Целью работ является изучение способов представления графов и практическая реализация алгоритмов для их генерации, анализа и преобразования.

Разработанная консольная программа на языке C++ объединяет результаты пяти лабораторных работ. Программа выполняет случайную генерацию графа и весовых матриц, вычисляет эксцентриситеты вершин, выделяет центр графа и диаметральные вершины, ищет маршруты, кратчайшие пути, точки сочленения и потоки, строит минимальный остов, кодирует его кодом Прюфера, находит максимальное независимое множество рёбер, строит эйлеров цикл и фундаментальную систему разрезов. Исходный граф, сформированный в первой работе, затем используется в последующих работах.

В отчёте рассмотрены математическое описание использованных алгоритмов, особенности их программной реализации и результаты работы программы на сгенерированных графах.
